%------------------------------------------------------------------------------
% Beamer presentation generated from:
%   1. introduction.md
%   2. chapter2.md
%   3. chapter3.md
%
% Build with:  pdflatex presentation.tex   (run twice for the outline/TOC)
%
% Figures are expected in ./Figures/ . If a figure file is absent, a labelled
% placeholder box is typeset instead so the document still compiles.
%------------------------------------------------------------------------------
\documentclass[10pt,aspectratio=169]{beamer}

\usetheme{Madrid}
\usecolortheme{dolphin}
\setbeamertemplate{navigation symbols}{}
\setbeamertemplate{caption}[numbered]

\usepackage[T1]{fontenc}
\usepackage[utf8]{inputenc}
\usepackage{amsmath,amssymb}
\usepackage{graphicx}
\usepackage{booktabs}
\usepackage{listings}

\graphicspath{{Figures/}{./}}

\lstset{
  basicstyle=\scriptsize\ttfamily,
  breaklines=true,
  columns=fullflexible,
  frame=single,
  showstringspaces=false,
  keepspaces=true,
  language=R,
  commentstyle=\color{gray},
  keywordstyle=\color{blue!60!black},
  stringstyle=\color{red!50!black}
}

% Include a figure, or a visible placeholder if the file is not present.
% #1 = path, #2 = caption / description
\newcommand{\figorbox}[2]{%
  \begin{center}
  \IfFileExists{#1}{%
    \includegraphics[width=0.85\linewidth,height=0.62\textheight,keepaspectratio]{#1}%
  }{%
    \fbox{\parbox[c][0.45\textheight][c]{0.8\linewidth}{\centering\itshape
      Figure not available at compile time:\\[0.3em]
      \texttt{\detokenize{#1}}\\[0.8em]
      #2}}%
  }%
  \end{center}%
}

\title{Statistics}
\subtitle{General Information, Introduction and Descriptive Statistics}
\date{\today}

\begin{document}

\begin{frame}
  \titlepage
\end{frame}

\begin{frame}{Outline}
  \tableofcontents
\end{frame}

%==============================================================================
\section{General Information}
%==============================================================================

\begin{frame}{Syllabus (1/2)}
  \begin{enumerate}
    \item Introduction; organizing and displaying data; summary measures.
    \item Probability distributions; discrete distributions: binomial, Poisson,
          geometric and negative binomial distributions.
    \item Continuous distributions: exponential, gamma, normal and Weibull
          distributions; QQ-plot.
    \item Sampling distribution; estimation and confidence intervals:
          Central Limit Theorem; parameter estimation (method of moments);
          confidence interval for mean and variance; comparison of two means and
          variances; introduction to hypothesis testing; Chi-square tests.
  \end{enumerate}
\end{frame}

\begin{frame}{Syllabus (2/2)}
  \begin{enumerate}
    \setcounter{enumi}{4}
    \item Regression and analysis of variance: simple regression and basic
          ANOVA tests.
    \item Practical statistics: the use of Minitab to analyse data and to
          interpret the results.
    \item Design of technical posters.
  \end{enumerate}
\end{frame}

\begin{frame}{Learning Outcomes (1/2)}
  \begin{itemize}
    \item Define a number of common distributions, fit them to data and test
          the goodness of fit.
    \item State and apply the Central Limit Theorem.
    \item Derive confidence intervals for the mean and variance of a sample, and
          compare the means and variances of two samples using appropriate
          hypothesis tests.
  \end{itemize}
\end{frame}

\begin{frame}{Learning Outcomes (2/2)}
  \begin{itemize}
    \item Conduct simple regression analysis and basic ANOVA tests, and assess
          the applicability of the models used to the data.
    \item Interpret the results of the statistical techniques taught on this
          module and present them accurately, succinctly and attractively to
          both specialists and non-specialists.
    \item Display technical material on a poster.
  \end{itemize}
\end{frame}

\begin{frame}{Assessment}
  \begin{itemize}
    \item \textbf{70\%} Closed book examination (2 hours).
    \item \textbf{30\%} Coursework, in the form of a poster and a short
          technical report.
    \item The pass mark is \textbf{50\%}.
  \end{itemize}
\end{frame}

\begin{frame}{References}
  \begin{itemize}
    \item Prem S. Mann, \emph{Introductory Statistics}, sixth edition.
    \item Anderson, Sweeney, Williams, Freeman, Shoesmith,
          \emph{Statistics for Business and Economics}.
    \item Walpole and Mayers,
          \emph{Probability and Statistics for Engineers and Scientists}.
    \item Hoel, \emph{Introduction to Mathematical Statistics}.
  \end{itemize}
\end{frame}

\begin{frame}{Software}
  Exercises may be completed in either \textbf{R} or via packages in
  \textbf{Python}; other statistical tools can also be used.

  \vspace{1em}
  If you are using something other than Python or R for the coursework, please
  confirm with the course lecturer before starting the work.
\end{frame}

%==============================================================================
\section{Introduction to Statistics}
%==============================================================================

\begin{frame}{Two Main Areas of Statistics}
  \begin{itemize}
    \item \emph{Mathematical statistics} develops and proves statistical
          theorems, rules and laws.
    \item \emph{Applied statistics} uses these to solve real-world problems.
  \end{itemize}

  \vspace{0.8em}
  We will deal mainly with applied statistics, which splits into two main areas:
  \begin{itemize}
    \item \emph{Descriptive statistics}: methods for organizing, displaying and
          describing data using tables, graphs and summary measures (mean, mode,
          median, etc.).
    \item \emph{Inferential statistics}: methods that use sample results to make
          decisions about a population.
  \end{itemize}
\end{frame}

\begin{frame}{Data: Elements, Variables and Observations}
  The starting point for statistical analysis is \emph{data} --- the plural of
  ``datum'', a piece of information. All the data collected in a particular
  study are referred to as a \emph{data set}.

  \vspace{0.6em}
  \begin{itemize}
    \item \emph{Elements} are the entities on which data are collected.
    \item A \emph{variable} is a characteristic of interest for the elements.
    \item The result of a measurement obtained for a particular element is an
          \emph{observed value} or \emph{observation}.
  \end{itemize}

  \vspace{0.6em}
  Measurements collected on each variable for every element in a study give a
  data set.
\end{frame}

\begin{frame}{Example: Data for Two Students}
  \begin{center}
    \begin{tabular}{|c|ccc|l|}
      \hline
              & \multicolumn{3}{c|}{Exam} &                 \\ \hline
      Name    & 1  & 2  & 3  & \textbf{Grade}               \\ \hline
      David   & 90 & 95 & 92 & A+                           \\
      Mary    & 85 & 90 & 91 & A                            \\ \hline
    \end{tabular}
  \end{center}

  \vspace{0.8em}
  \begin{itemize}
    \item The two \emph{elements} are ``David'' and ``Mary''.
    \item The \emph{variables} are ``Exam 1'', ``Exam 2'', ``Exam 3'' and
          ``Grade''.
    \item For the Exam variables the observations are \emph{numerical}; for the
          Grade variable they are \emph{categorical}.
  \end{itemize}
\end{frame}

\begin{frame}{Qualitative and Quantitative Data}
  \begin{itemize}
    \item \emph{Qualitative} (categorical) \emph{data} include labels or names
          used to identify an attribute of each element.
    \item \emph{Quantitative data} require numerical values indicating how much
          or how many. They can be
          \begin{itemize}
            \item \textbf{discrete} (number of houses, cars, \dots), or
            \item \textbf{continuous} (height, weight, \dots).
          \end{itemize}
  \end{itemize}

  \vspace{0.6em}
  A \emph{qualitative variable} is a variable with qualitative data; a
  \emph{quantitative variable} is a variable with quantitative data.
\end{frame}

\begin{frame}{Why the Type of Data Matters}
  An appropriate statistical analysis depends on whether the variable is
  quantitative or qualitative.

  \vspace{0.8em}
  The statistical analysis for \textbf{qualitative} data is rather limited:
  \begin{itemize}
    \item determine frequencies and relative frequencies;
    \item represent these using a bar chart or a pie chart.
  \end{itemize}

  \vspace{0.8em}
  There are many more alternatives for the statistical analysis of
  \textbf{quantitative} data.
\end{frame}

\begin{frame}{Cross-Sectional and Time Series Data}
  \begin{itemize}
    \item \emph{Cross-sectional data} are collected on different elements at
          approximately the same point in time.
    \item \emph{Time series data} are collected over several time periods on the
          same element(s).
  \end{itemize}
\end{frame}

\begin{frame}{Data Sources}
  Data can be obtained from existing sources, surveys, a census, and statistical
  studies designed to collect new data.

  \vspace{0.8em}
  Many situations require data for a large group of elements (individuals,
  companies, voters, households, products, customers, \dots). Because of time,
  cost or other reasons, data are often collected from only a small portion of
  the group.
\end{frame}

\begin{frame}{Population, Sample, Census}
  \begin{itemize}
    \item The group of all elements of interest in a particular study is the
          \emph{population}.
    \item The smaller group of the population selected for the study is the
          \emph{sample}.
    \item A \emph{census} is the process of conducting a survey to collect data
          on the entire population.
    \item A \emph{sample survey} collects data from a sample.
  \end{itemize}

  \vspace{0.8em}
  The methods of statistical inference use sample results to make estimates or
  test hypotheses about the characteristics of a population.
\end{frame}

\begin{frame}{Where We Go Next}
  We start by defining some descriptive statistics that are typically collected
  as part of an \emph{exploratory data analysis}, and which can detect
  ``global'' features of data such as shape, peaks and outliers.

  \vspace{1em}
  We then move on to more concrete distributions of data, and introduce
  inferential statistics.
\end{frame}

%==============================================================================
\section{Descriptive Statistics}
%==============================================================================

\subsection{Graphical Data Presentations}

\begin{frame}{Graphical Data Presentations}
  Even a small set of data is often difficult to interpret straight from its
  values.

  \vspace{1em}
  These so-called \emph{raw data} are organised into tables and displayed using
  plots and other graphics.
\end{frame}

\begin{frame}{Summarising Qualitative Data}
  For a \textbf{qualitative} data set we can construct \emph{frequency} and
  \emph{relative frequency} distributions.

  \begin{itemize}
    \item A \emph{frequency distribution} is a tabular summary showing the
          number (frequency) of items in each of several non-overlapping
          classes. Often displayed using a \emph{bar chart}.
    \item A \emph{relative frequency distribution} displays the frequency of
          each class divided by the number of elements in the whole data set.
          Often displayed using a \emph{pie chart}.
  \end{itemize}

  With $n$ observations,
  \[
    \text{Relative frequency of class} = \frac{\text{Frequency of class}}{n}.
  \]
\end{frame}

\begin{frame}{Example: Last Names}
  \begin{block}{Data set}
    This example data set includes 50 individuals with one of the six most
    common last names in Ireland. The data file \texttt{names.txt} is on
    Blackboard.
  \end{block}

  \vspace{0.8em}
  We can draw a bar chart to show the frequencies of last names within the
  data set.
\end{frame}

\begin{frame}[fragile]{Example: Last Names --- R Code}
\begin{lstlisting}
# Read in data
names <- readLines("names.txt")
field_name <- "last_names"
names_df <- data.frame(name = names[-1], stringsAsFactors = FALSE)
colnames(names_df) <- field_name

# Find frequencies
frequency_df <- as.data.frame(table(names_df[[1]]))
colnames(frequency_df) <- c(field_name, "Frequency")

# Draw a bar plot with lime green coloured bars
ggplot(frequency_df, aes(x = last_names, y = Frequency)) +
  geom_bar(stat = "identity", fill = "limegreen") +
  labs(
    title = "Frequency of Last Names",
    x = names[1],
    y = "Frequency"
  )
\end{lstlisting}
\end{frame}

\begin{frame}{Example: Last Names --- Bar Chart}
  \figorbox{Figures/names_bar.png}{Bar chart of last names}
\end{frame}

\begin{frame}{Classes for Quantitative Data}
  With \textbf{quantitative} data we must be more careful in defining classes.

  \begin{itemize}
    \item A class is defined by its upper and lower limits: a data point belongs
          to a class if its value lies between those limits.
    \item There should be no overlap between classes.
    \item In the majority of examples, using the same width for each class works
          best.
    \item The number of classes depends on the number of data points and their
          range; typically between 5 and 20 classes are used.
  \end{itemize}

  Choosing the number of classes requires judgment --- a decision that best
  displays any natural grouping and variation in the data.

  \vspace{0.5em}
  The relative frequency is defined as before and represents the proportion of
  observations belonging to a specific class.
\end{frame}

\begin{frame}{Example: Computer Usage --- Histograms}
  \begin{block}{Data set}
    Data on personal computer usage by persons aged 12 and older during one
    week, for a sample of 50 persons, in hours. The data are in
    \texttt{computer.txt} on Blackboard.
  \end{block}

  A \emph{histogram} is a graphical representation of quantitative data showing
  the frequency or relative frequency of data values within classes. It is
  useful for understanding the shape or \emph{skewness} of a numerical variable.

  \vspace{0.5em}
  \textbf{How does a histogram differ from a bar chart?}
  A histogram displays numerical data, whereas a bar chart is typically used for
  categorical data.

  \vspace{0.5em}
  The frequency of each class is shown by a rectangle whose base is determined
  by the class limits on the horizontal axis, and whose height is the
  corresponding frequency or relative frequency.
\end{frame}

\begin{frame}[fragile]{Example: Computer Usage --- R Code}
\begin{lstlisting}
# Read in computer data
computer <- readLines("computer.txt")
computer_field_name <- "Hours"
computer_df <- data.frame(computer = computer[-1])
colnames(computer_df) <- computer_field_name
computer_df$Hours <- as.double(computer_df$Hours)

# Draw histogram
ggplot(computer_df, aes(x = Hours)) +
  geom_histogram(
    bins = 12,
    fill = "orange"
  ) +
  labs(
    title = "Histogram of hours spent on a computer",
    x = "Hours",
    y = "Frequency"
  )
\end{lstlisting}
\end{frame}

\begin{frame}{Example: Computer Usage --- Histogram}
  \figorbox{Figures/HistogramofHours.png}{Histogram of computer hours}
\end{frame}

\begin{frame}{Cumulative Frequency Distribution}
  \begin{itemize}
    \item The \emph{cumulative frequency distribution} shows the number of data
          items with values less than or equal to the upper class limit of each
          class.
    \item The \emph{cumulative relative frequency distribution} shows the
          proportion of data items with values less than or equal to the upper
          limit of each class.
    \item A graph of a cumulative distribution is called an \emph{ogive}: data
          values on the horizontal axis, cumulative frequencies or cumulative
          relative frequencies on the vertical axis.
  \end{itemize}

  \vspace{0.5em}
  We return to this later in the chapter when we introduce the cumulative
  distribution function.
\end{frame}

\begin{frame}[fragile]{Empirical CDF --- R Code}
  For example, we can plot the empirical cumulative distribution function (ECDF)
  for the computer example.
\begin{lstlisting}
ggplot(computer_df, aes(x = Hours)) +
  stat_ecdf() +
  labs(
    title = "Cumulative Distribution Function of Hours",
    x = "Hours",
    y = "Cumulative Probability"
  )
\end{lstlisting}
\end{frame}

\begin{frame}{Empirical CDF of Computer Hours}
  \figorbox{Figures/EmpiricalCDFofComputerHours.png}{ECDF of computer hours}
\end{frame}

\begin{frame}{Scatter Diagram and Trend Line}
  A \emph{scatter diagram} is a graphical representation of the relationship
  between two quantitative variables. A \emph{trend line} is a line that gives
  an approximation of that relationship.

  \begin{block}{Example: Computer Store}
    The data have three columns: the week number, the number of commercials
    shown, and the volume of sales at the equipment store that week. Available
    on Blackboard as \texttt{store.txt}.
  \end{block}

  We can draw a scatterplot to illustrate the relationship between the number of
  commercials and sales volume.
\end{frame}

\begin{frame}[fragile]{Computer Store --- R Code}
\begin{lstlisting}
# Read in equipment store data
equipment_df <- read.table(
  "store.txt",
  header = TRUE,
  sep = "\t"
)

ggplot(equipment_df, aes(x = `No..of.Commercials`, y = `Sales.Volume`)) +
  geom_point() +
  geom_smooth(method = "lm", se = TRUE) +
  labs(
    title = "Commercials vs Sales Volume",
    x = "Number of Commercials",
    y = "Sales Volume"
  )
\end{lstlisting}
\end{frame}

\begin{frame}{Computer Store --- Scatter Plot}
  \figorbox{Figures/scatterStore.png}{Scatter plot of equipment store data}
\end{frame}

\subsection{Numerical Descriptive Measures}

\begin{frame}{Numerical Descriptive Measures}
  Numerical descriptive measures are measures of \textbf{location},
  \textbf{dispersion}, \textbf{shape} and \textbf{association}.

  \begin{itemize}
    \item Calculated for sample data, they are called \emph{sample statistics}.
    \item Calculated for population data, they are called
          \emph{population parameters}.
    \item A sample statistic is the \emph{point estimator} of the corresponding
          population parameter.
  \end{itemize}

  \begin{block}{Notation convention}
    \begin{itemize}
      \item $X$ (capital) denotes a variable, whose observations can be made.
      \item $x$ (small) denotes a generic value that the variable $X$ may take.
    \end{itemize}
    For example, $X$ could be the height of people, and $x = 185\,\mathrm{cm}$ is
    a value of $X$ for a particular person.
  \end{block}
\end{frame}

\begin{frame}{Measures of Location: The Mean}
  The \emph{mean} (average value) is the most frequently used measure of central
  location.

  \begin{itemize}
    \item The \textbf{sample mean} of $\mathbf{x} = (x_1, x_2, \ldots, x_n)$ is
          denoted $\bar{x}$.
    \item The \textbf{population mean} of
          $\mathbf{x} = (x_1, x_2, \ldots, x_N)$ is denoted $\mu$.
  \end{itemize}

  \[
    \bar{x} = \frac{\sum_{i=1}^{n} x_i}{n}
    \qquad\text{and}\qquad
    \mu = \frac{\sum_{i=1}^{N} x_i}{N},
  \]
  where $n$ is the sample size, $N$ the population size, and $x_i$ the value of
  the variable $X$ on the $i$-th observation.

  \vspace{0.5em}
  The expression for the mean includes all values of the data set and, as a
  result, it is \textbf{sensitive to outliers}.
\end{frame}

\begin{frame}{Measures of Location: Median and Mode}
  \begin{block}{Median}
    The value of the ``middle'' term when the data are arranged in increasing
    order. If $n$ is odd it is the middle value; if $n$ is even it is the average
    of the two middle values. The median is \textbf{not influenced by outliers}.
  \end{block}

  \begin{block}{Mode}
    The value in a data set that occurs with the highest frequency. It can be
    determined for both quantitative and qualitative data.

    A data set can have no mode (all values equally frequent), one mode, or
    several modes. The modes of multimodal data are not very helpful.
  \end{block}
\end{frame}

\begin{frame}{Example: Trades of Lloyds Bank Group Shares}
  The last 5 trades of Lloyds Bank Group shares on 22 September 2017
  (data from the London Stock Exchange).

  \vspace{0.5em}
  \begin{center}
    \begin{tabular}{llrr}
      \toprule
      Time & Price & Volume & Trade value \\
      \midrule
      17:11:37 & 66.98 & 93,163 & 62,400.53 \\
      17:11:29 & 67.03 & 80,204 & 53,760.78 \\
      17:11:29 & 67.03 & 13,178 &  8,833.22 \\
      17:03:06 & 67.15 & 70,800 & 47,542.20 \\
      17:02:36 & 67.00 &  9,064 &  6,072.64 \\
      \bottomrule
    \end{tabular}
  \end{center}
\end{frame}

\begin{frame}{Lloyds Example: Mean, Median and Mode}
  The \textbf{mean} price is
  \[
    \bar{p} = \frac{66.98 + 67.03 + 67.03 + 67.15 + 67.00}{5}
            = 67.038 \approx 67.04 \ \text{(2 d.p.)}.
  \]

  Here $n = 5$, so the \textbf{median} is the third highest price. The ordered
  list is $(66.98,\ 67.00,\ 67.03,\ 67.03,\ 67.15)$, and the third value is
  $67.03$.

  The price $67.03$ appears twice, so this is the \textbf{mode}. For continuous
  numerical data it may be more appropriate to find the range of values that
  contains the most data points.

  \vspace{0.8em}
  Deciding which measure to use (mean, median or mode) depends on your data set.
\end{frame}

\begin{frame}{Measures of Variability}
  The most common measures of variability are the \textbf{range},
  \textbf{interquartile range}, \textbf{variance} and
  \textbf{standard deviation}.

  \begin{itemize}
    \item The \emph{range} is the difference between the largest and smallest
          values. It is based on only two observations and is highly influenced
          by outliers.
    \item The \emph{interquartile range} is the range of the middle 50\% of the
          data, and so does not depend on outliers.
    \item The \emph{variance} takes into account all of the data. It is based on
          the \emph{deviation about the mean},
          \[
            x_i - \bar{x}.
          \]
  \end{itemize}
\end{frame}

\begin{frame}{Variance and Standard Deviation}
  Population variance $\sigma^2$, for a population of size $N$ with mean $\mu$:
  \[
    \sigma^2 = \frac{\sum_{i=1}^{N} (x_i - \mu)^2}{N-1}.
  \]
  Sample variance $s^2$, for a sample of size $n$ with mean $\bar{x}$:
  \[
    s^2 = \frac{\sum_{i=1}^{n} (x_i - \bar{x})^2}{n-1}.
  \]
  The statistic $s^2$ gives an unbiased estimate of the population variance.
  Note the divisor $(n-1)$, not $n$ as one might expect.

  \vspace{0.5em}
  The variance is useful when comparing variability between two or more
  variables, samples or populations: the one with the largest variance shows the
  greatest variability.

  \vspace{0.5em}
  The \emph{standard deviation} is the positive square root of the variance. It
  is measured in the same units as the original data, and is consequently common
  in practice.
\end{frame}

\begin{frame}{Coefficient of Variation}
  The \emph{coefficient of variation} measures how large the standard deviation
  is relative to the mean:
  \[
    \frac{\text{Standard deviation}}{\text{Mean}} \times 100\%.
  \]

  \vspace{1em}
  It is a descriptive statistic that indicates the size of the standard
  deviation relative to the mean, and is independent of the units of the data.
\end{frame}

\begin{frame}{Percentiles}
  A \emph{percentile} provides information on how the ordered data are spread
  over the range of the data. The $p$-th percentile is a value such that at
  least $p$ percent of the data are less than or equal to it, and at least
  $(100-p)$ percent are greater than or equal to it.

  \begin{block}{Calculating percentiles for a sample $x_1, x_2, \ldots, x_n$}
    \begin{enumerate}
      \item Arrange the data in increasing order,
            $x_{(1)} \le x_{(2)} \le \cdots \le x_{(n)}$.
      \item Calculate $i = \left(\frac{p}{100}\right) n$, where $p$ is the
            percentile of interest and $n$ the sample size.
      \item If $i$ is \textbf{not} an integer, round up; the resulting value of
            $i$ gives the position of the $p$-th percentile. If $i$ \textbf{is}
            an integer, the $p$-th percentile is the average of the values in
            positions $i$ and $i+1$.
    \end{enumerate}
  \end{block}
\end{frame}

\begin{frame}{Example: Waiting Time}
  \begin{block}{Data set}
    The waiting time data is a list of 24 values describing the time spent
    waiting for a service. They can be found in the file
    \texttt{WaitingTime.txt}.
  \end{block}

  \textbf{(a) Compute the 30th percentile.} Here $n = 24$ and $p = 30$:
  \[
    i = \left(\frac{p}{100}\right) \times n = \frac{30}{100} \times 24 = 7.2 .
  \]
  Hence the 30th percentile is
  \[
    x_{(\lceil i \rceil)} = x_{(\lceil 7.2 \rceil)} = x_{(8)} = 30,
  \]
  where $\lceil i \rceil$ is the smallest integer greater than or equal to $i$.
\end{frame}

\begin{frame}{Example: Waiting Time --- Quartiles}
  \textbf{(b) Compute the 25th percentile.} Here $n = 24$, $p = 25$:
  \[
    i = \frac{25}{100} \times 24 = 6 \quad (\text{an integer}),
  \]
  \[
    Q_1 = \frac{x_{(i)} + x_{(i+1)}}{2}
        = \frac{x_{(6)} + x_{(7)}}{2}
        = \frac{26 + 28}{2} = 27 .
  \]

  \textbf{(c) Compute the 75th percentile.} Here $n = 24$, $p = 75$:
  \[
    i = \frac{75}{100} \times 24 = 18 \quad (\text{an integer}),
  \]
  \[
    Q_3 = \frac{x_{(18)} + x_{(19)}}{2} = \frac{61 + 62}{2} = 61.5 .
  \]
\end{frame}

\begin{frame}{Quartiles}
  The \emph{quartiles} $Q_1$, $Q_2$ and $Q_3$ divide an ordered data set into
  four equal parts.

  \begin{itemize}
    \item $Q_1$ is the 25th percentile: 25\% of the data set is equal to or
          smaller than $Q_1$.
    \item $Q_2$ is equal to the median.
    \item $Q_3$ is the 75th percentile.
    \item The \emph{interquartile range} is $\mathrm{IQR} = Q_3 - Q_1$.
  \end{itemize}

  \begin{block}{Steps to calculate the quartiles}
    \begin{enumerate}
      \item Arrange the observations in non-decreasing order and locate the
            median $M$.
      \item $Q_1$ is the median of the observations whose positions in the
            ordered list lie to the left of $M$.
      \item $Q_3$ is the median of the observations whose positions in the
            ordered list lie to the right of $M$.
    \end{enumerate}
  \end{block}
\end{frame}

\begin{frame}{The Five-Number Summary and Boxplots}
  The five-number summary of a set of observations consists of
  \[
    \text{Minimum}, \quad Q_1, \quad M, \quad Q_3, \quad \text{Maximum}.
  \]

  These are well represented in a \emph{boxplot}:
  \begin{itemize}
    \item a central box spanning the quartiles $Q_1$ and $Q_3$;
    \item a line in the box marking the median $M$;
    \item lines extending from the box out to the lower and upper limits;
    \item the $1.5 \times \mathrm{IQR}$ rule determines those limits:
      \[
        \begin{aligned}
          \text{lower limit} &= M - 1.5 \times \mathrm{IQR} \\
          \text{upper limit} &= M + 1.5 \times \mathrm{IQR}
        \end{aligned}
      \]
    \item all other observations beyond the limits are marked as whiskers; they
          are suspected outliers.
  \end{itemize}
\end{frame}

\begin{frame}{Boxplot: Waiting Time Example}
  We have already calculated
  \[
    Q_1 = 27, \quad M = 45.5, \quad Q_3 = 61.5,
  \]
  so
  \[
    \mathrm{IQR} = Q_3 - Q_1 = 61.5 - 27 = 34.5 .
  \]
  The lower and upper limits are
  \[
    L = 45.5 - 1.5 \times 34.5 = -6.25,
    \qquad
    U = 45.5 + 1.5 \times 34.5 = 117.15 .
  \]
\end{frame}

\begin{frame}{Boxplot for the Waiting Time Data}
  \figorbox{Figures/BoxPlotWaiting.png}{Boxplot for the waiting time data set}
\end{frame}

\subsection{Distributions and Shape}

\begin{frame}[fragile]{Probability Density Functions --- R Code}
  A probability density function $f(x)$, or PDF (defined properly later),
  provides a more complete description of how the data are distributed. For
  example, we might fit a distribution to the waiting time data.
\begin{lstlisting}
ggplot(waiting_df, aes(x = Waiting.Time)) +
  geom_histogram(
    aes(y = after_stat(density)),
    bins = 13,
    fill = "limegreen",
  ) +
  stat_function(
    fun = dnorm,
    args = list(
      mean = mean(waiting_df$Waiting.Time),
      sd = sd(waiting_df$Waiting.Time)
    ),
    colour = "blue",
    linewidth = 1
  ) +
  labs(x = "Waiting times", y = "Density")
\end{lstlisting}
\end{frame}

\begin{frame}{Waiting Times with a Fitted Normal Distribution}
  \figorbox{Figures/WaitingHistogramNormal.png}%
    {Histogram of the waiting times data with a fitted normal distribution}
\end{frame}

\begin{frame}{The Normal Distribution}
  The fitted curve is completely determined by the sample mean
  ($\bar{x} = 47.04$) and the standard deviation ($s = 31.87$) of the data.

  The density function takes the sample mean as its true mean $\mu$ and the
  sample standard deviation as its true standard deviation $\sigma$:
  \[
    f(x) = \frac{1}{\sqrt{2\pi\sigma^2}}\,
           e^{-\frac{(x-\mu)^2}{2\sigma^2}} .
  \]

  This is the \textbf{normal} (or Gaussian) distribution, often denoted
  $N(\mu, \sigma)$. Its location and shape are completely determined by
  $(\mu, \sigma)$, and the density is symmetric about $\mu$.

  \vspace{0.5em}
  It is a density function because
  \begin{enumerate}
    \item $f(x) \ge 0$ for all $x \in \Re$;
    \item the area underneath the curve is 1:
          $\displaystyle\int_{-\infty}^{+\infty} f(x)\,dx = 1$.
  \end{enumerate}
\end{frame}

\begin{frame}{The 68--95--99.7 Rule}
  Although different $(\mu, \sigma)$ generate different curves, normal
  distributions share a common feature.

  \vspace{1em}
  In the normal distribution with mean $\mu$ and standard deviation $\sigma$:
  \begin{itemize}
    \item approximately $68\%$ of observations fall within $\sigma$ of the mean
          $\mu$;
    \item approximately $95\%$ fall within $2\sigma$ of the mean $\mu$;
    \item approximately $99.7\%$ fall within $3\sigma$ of the mean $\mu$.
  \end{itemize}
\end{frame}

\begin{frame}{The Standard Normal Distribution}
  Introduce a new variable $Z$ such that
  \[
    Z = \frac{X - \mu}{\sigma},
    \qquad\text{equivalently}\qquad
    X = \sigma Z + \mu .
  \]
  It follows from $\int f(x)\,dx = 1$ over $(-\infty, \infty)$ that
  \[
    \begin{aligned}
      1 &= \int_{-\infty}^{+\infty} f(x)\,dx \\
        &= \int_{-\infty}^{+\infty}
             \frac{1}{\sigma\sqrt{2\pi}}
             e^{-\frac{1}{2}\left(\frac{x-\mu}{\sigma}\right)^2} dx \\
        &= \int_{-\infty}^{+\infty}
             \frac{1}{\sigma\sqrt{2\pi}} e^{-\frac{1}{2}z^2} \sigma\,dz \\
        &= \int_{-\infty}^{+\infty}
             \frac{1}{\sqrt{2\pi}} e^{-\frac{1}{2}z^2}\,dz .
    \end{aligned}
  \]
\end{frame}

\begin{frame}{The Standard Normal Density}
  The function
  \[
    \phi(z) = \frac{1}{\sqrt{2\pi}}\, e^{-\frac{1}{2}z^2}
  \]
  is a density function corresponding to a normal distribution with mean $0$ and
  standard deviation $\sigma = 1$.

  \vspace{1em}
  Its distribution is denoted $N(0,1)$ and is known as the
  \textbf{standard normal distribution}.
\end{frame}

\begin{frame}{Cumulative Distribution Function (CDF)}
  Let $X$ follow a distribution with pdf $f(x)$. The cumulative distribution at
  a given value $x$ is
  \[
    F(x) = \int_{-\infty}^{x} f(t)\,dt .
  \]

  The CDF of the standard normal distribution is frequently used to estimate
  confidence intervals on estimates; we will see later why.

  \vspace{0.5em}
  Let $Z \sim N(0,1)$ with density $\phi(z)$. Its cumulative distribution is
  \[
    \Phi(z) = \int_{-\infty}^{z} \phi(t)\,dt .
  \]
  It follows from the $Z$ transformation that
  \[
    F(x) = \Phi\!\left( \frac{x-\mu}{\sigma} \right),
  \]
  so $F(x)$ can be obtained from $\Phi(z)$ with $z = (x-\mu)/\sigma$. Tables of
  $\Phi(z)$ are widely available.
\end{frame}

\begin{frame}{Relationship Between $F(x)$ and $\Phi(z)$}
  \figorbox{Figures/NormalDistributionZTransformation.jpg}%
    {Relationship between $F(x)$ and $\Phi(z)$ (source: Internet)}
\end{frame}

\begin{frame}{Skewness}
  The normal distribution has a symmetric shape, but many shapes are
  non-symmetric --- a shape may be skewed to the left or to the right.

  The measure of skewness is
  \[
    \texttt{skewness} = \frac{n}{(n-1)(n-2)}
      \sum_{i=1}^{n} \left( \frac{x_i - \bar{x}}{s} \right)^{3} .
  \]
  Skewness can be easily calculated using statistical software.

  \begin{itemize}
    \item A distribution is \textbf{skewed to the left} (negatively skewed) if
          its tail extends further to the left.
    \item A distribution is \textbf{right skewed} if its tail extends further to
          the right.
    \item A distribution is \textbf{symmetrical} if there is no obvious skewness.
  \end{itemize}
\end{frame}

\begin{frame}{Examples of Different Values of Skewness}
  \figorbox{Figures/skewness.jpeg}{Examples of different values of skewness}
\end{frame}

\begin{frame}{Z-Scores}
  The \emph{z-scores} give the relative location of data items within a data
  set. In a sample of $n$ observations $x_1, x_2, \ldots, x_n$, the z-score
  $z_i$ for observation $x_i$ is
  \[
    z_i = \frac{x_i - \bar{x}}{s}, \qquad i = 1, 2, \ldots, n,
  \]
  where $\bar{x}$ is the sample mean and $s$ the sample standard deviation. Note
  the similarity to the $Z$ transformation.

  The z-score is often called the \emph{standardized value} and is interpreted
  as the number of standard deviations $x_i$ lies from the mean $\bar{x}$. For
  example, $z_1 = 1.7$ indicates that $x_1$ is 1.7 standard deviations greater
  than the mean; $z_1 = -0.5$ indicates it is $0.5$ standard deviations smaller.
\end{frame}

\begin{frame}{Outliers}
  Extreme values (unusually large or unusually small) of a data set are called
  \emph{outliers}.

  \vspace{1em}
  Identification and review of outliers is important.

  \vspace{1em}
  Typically, any data value with a z-score less than $-3$ or greater than $+3$
  is treated as an outlier.
\end{frame}

\subsection{Measures of Association}

\begin{frame}{Measures of Association}
  We are often interested in the relationship between two or more variables.

  Returning to the equipment store example and its scatter plot: the plot
  suggests a \textbf{positive relationship} between the number of commercials
  shown and the volume of sales, and that a straight line could approximate the
  relationship.

  \begin{itemize}
    \item A \textbf{positive} relationship: as the number of commercials
          increases, the number of sales also increases.
    \item A \textbf{negative} relationship: as the number of commercials
          increases, the number of sales decreases.
  \end{itemize}
\end{frame}

\begin{frame}{Covariance}
  The \emph{covariance} is a measure of the linear relationship between two
  variables. For a sample of size $n$ with observations
  $(x_1, y_1), (x_2, y_2), \ldots, (x_n, y_n)$, the sample covariance $s_{xy}$
  is
  \[
    s_{xy} = \frac{\sum_{i=1}^{n} (x_i - \bar{x})(y_i - \bar{y})}{n-1},
  \]
  and the covariance of a population of size $N$ is
  \[
    \sigma_{xy} = \frac{\sum_{i=1}^{N} (x_i - \mu_x)(y_i - \mu_y)}{N-1}.
  \]

  \begin{alertblock}{Exercise}
    Calculate the covariance $s_{xy}$ for the previous example.
    \hfill (Answer: 11)
  \end{alertblock}
\end{frame}

\begin{frame}{The Correlation Coefficient}
  A problem with covariance as a measure of the strength of a linear
  relationship is that its value depends on the units of measurement. We need a
  measure independent of those units.

  The \textbf{correlation coefficient} $r_{xy}$ is
  \[
    r_{xy} = \frac{s_{xy}}{s_x s_y},
  \]
  where $s_{xy}$ is the sample covariance and $s_x$, $s_y$ are the sample
  standard deviations of $x$ and $y$.

  \begin{alertblock}{Exercise}
    Calculate the sample correlation coefficient for the example above. Does
    this help with determining the strength of the relationship?
    \hfill (Answer: 0.93)
  \end{alertblock}
\end{frame}

\begin{frame}{Population Correlation and Its Range}
  For population data,
  \[
    \rho_{xy} = \frac{\sigma_{xy}}{\sigma_x \sigma_y},
  \]
  where $\rho_{xy}$ is the population correlation coefficient, $\sigma_{xy}$ the
  population covariance, and $\sigma_x$, $\sigma_y$ the population standard
  deviations.

  The correlation has a limited range, $-1 \le \rho_{xy} \le 1$:
  \begin{itemize}
    \item $\rho_{xy} = 1$ indicates a perfect positive linear relationship;
    \item $\rho_{xy} = -1$ indicates a perfect negative relationship;
    \item $\rho_{xy} = 0$ indicates no relationship between $x$ and $y$.
  \end{itemize}

  \begin{alertblock}{Exercise}
    Let $X$ and $Y$ be random variables with finite means and variances. Using
    the expression for $\rho_{xy}$ above, show that $Y = a + bX$ if and only if
    $\rho = \pm 1$.
  \end{alertblock}
\end{frame}

\begin{frame}{The Weighted Mean}
  The \emph{weighted sample mean} attaches more importance (weight) to some
  observations than to others:
  \[
    \bar{x} = \frac{\sum_{i=1}^{n} w_i x_i}{\sum_{i=1}^{n} w_i},
  \]
  where $x_i$ is the value of observation $i$ and $w_i$ is the weight for
  observation $i$.
\end{frame}

\begin{frame}{Grouped Data}
  If data are grouped into $k$ classes (groups or bins) with mid values
  $M_1, M_2, \ldots, M_k$ and corresponding class frequencies
  $f_1, f_2, \ldots, f_k$, the mean is
  \[
    \bar{x} = \frac{\sum_{i=1}^{k} f_i M_i}{n},
  \]
  and the sample variance for grouped data is
  \[
    s^2 = \frac{\sum_{i=1}^{k} f_i (M_i - \bar{x})^2}{n-1},
  \]
  where $f_1 + f_2 + \cdots + f_k = n$.
\end{frame}

\begin{frame}{Summary}
  \begin{itemize}
    \item Descriptive statistics organise and display raw data: frequency
          tables, bar charts, pie charts, histograms, ogives and scatter
          diagrams.
    \item Numerical measures describe location (mean, median, mode), variability
          (range, IQR, variance, standard deviation, coefficient of variation),
          shape (skewness) and association (covariance, correlation).
    \item The normal distribution and the $Z$ transformation underpin much of
          what follows.
    \item Next: probability distributions and inferential statistics.
  \end{itemize}
\end{frame}

\end{document}
